\documentclass[11pt,letterpaper]{article}
\usepackage[margin=1in]{geometry}
\usepackage[T1]{fontenc}
\usepackage{lmodern}
\usepackage{microtype}
\usepackage{amsmath,amssymb,amsthm,mathtools}
\usepackage{xcolor}
\usepackage{enumitem}
\usepackage[hidelinks]{hyperref}
\hypersetup{pdftitle={A quantitative refinement of a depth-three circuit lower bound},pdfauthor={Tom Ayre},pdfsubject={Quantitative sparsification and Boolean circuit lower bounds}}
\newtheorem{theorem}{Theorem}[section]
\newtheorem{lemma}[theorem]{Lemma}
\newtheorem{proposition}[theorem]{Proposition}
\theoremstyle{definition}
\newtheorem{definition}[theorem]{Definition}
\numberwithin{equation}{section}
\newcommand{\E}{\mathbb E}
\newcommand{\Prb}{\mathbb P}
\newcommand{\bits}{\{0,1\}}
\DeclareMathOperator{\corr}{corr}
\setlist[itemize]{itemsep=3pt,topsep=5pt}
\setlength{\emergencystretch}{2em}
\clubpenalty=10000
\widowpenalty=10000
\title{A quantitative refinement of a\ depth-three circuit lower bound}
\author{Tom Ayre}
\date{October 6, 2026}
\begin{document}
\maketitle
\begin{abstract}
OpenAI's depth-three circuit construction gives a single polynomial-time
language whose $n$-bit membership function requires
$2^{\omega(\sqrt n)}$ total gates in unrestricted OR--AND--OR circuits.
We derive the explicit lower bound
$2^{c\sqrt{n\ln n/\ln\ln n}}$ for an absolute $c>0$, at every sufficiently
large input length, for the same language. The additional input is the
published quantitative form of disjoint CNF sparsification, with clause
density at most $(Kb)^{Kb}$ at width $b$. This permits the target width to
grow as $\Theta(\ln n/\ln\ln n)$. We verify the counting, moment, restriction,
and gate-count estimates uniformly for this schedule. We also reproduce
the source's finite restriction argument, including the backwards path
count that removes dependence on the number of clauses. The refinement
changes the proof parameters, not the language or its evaluation algorithm.
\end{abstract}

\section{Introduction}
For a Boolean function $f$, let $S_3(f)$ denote the minimum total number
of AND and OR gates in an exact OR--AND--OR circuit for $f$. Bottom gates
read input literals and optional constants; middle gates read bottom-gate
outputs; and the output is an OR of middle-gate outputs. Fan-in and
fan-out are unrestricted. Input negations are free, and all bottom,
middle, and output gates are counted. Unary gates permit functions of
smaller depth in this model.

OpenAI~\cite{OAI26} constructs one language $L\in\mathsf P$ such that its
$n$-bit membership function $f_n$ satisfies
$\log_2 S_3(f_n)/\sqrt n\to\infty$. Its argument first fixes a desired
constant in the circuit exponent, then chooses a constant CNF width, and
finally takes the input length sufficiently large. The present note
extracts an explicit rate by allowing that width to grow. The language is
recalled completely in Section~\ref{sec:language}.

\begin{theorem}\label{thm:main}
For the language $L$ of Definition~\ref{def:language}, there exist absolute
constants $c>0$ and $n_0$ such that
\begin{equation}\label{eq:main}
 S_3(f_n)>2^{c\sqrt{n\ln n/\ln\ln n}}
 \qquad\text{for every }n\ge n_0.
\end{equation}
The language is decided by one deterministic polynomial-time algorithm.
For any absolute $K\ge2$ satisfying Lemma~\ref{lem:sparsification},
one may take $c=1/(6144\sqrt K)$.
\end{theorem}

The displayed value of $c$ is explicit in terms of $K$; no numerical
value of the sparsification constant is claimed. The threshold $n_0$ is
absolute once this constant has been fixed.

\paragraph{Contribution and attribution.}
The language, clause-count-independent restriction lemma, and circuit
contradiction are from~\cite{OAI26}. The restriction proof reproduced in
Appendix~\ref{app:restriction} is an adaptation of that source, not a new
restriction lemma. The extra ingredient is the quantitative sparsification
estimate of Calabro, Impagliazzo, and Paturi~\cite{CIP}, combined with the
disjoint refinement of Dell et al.~\cite{DHMTW}; Williams~\cite{Williams}
explicitly uses this quantitative counting form. Our contribution is the
growing-parameter schedule and the uniform estimates that yield
\eqref{eq:main} for the existing construction. We regard this as a
potentially new quantitative refinement of the existing result; publication
priority remains unconfirmed.

The distinction between unrestricted and bounded-bottom-fan-in circuits
is essential. Recent exponential bounds for majority and inner product
concern restricted models~\cite{GKKPRST,GKPRT26IP}. The stronger unrestricted
majority bound $2^{\Omega(\sqrt{n\log n})}$ is presented as a conjectural
target in~\cite{GKKPRST}. The rate in~\eqref{eq:main} is not stated in
the source manuscript or its accompanying formalization. A targeted
literature search found no prior statement of this quantitative refinement;
this observation does not establish publication priority.

\paragraph{Proof overview.}
The construction has data bits and parameter bits. Suitable fixed parameter
bits give a data slice whose sign has small correlation with every narrow
CNF on most random restrictions. Sparsification makes the family of tests
small enough for a limited-independence moment bound. The restriction lemma
transfers this correlation bound to wider CNFs. After deleting very wide
bottom clauses from a hypothetical small circuit, each middle gate accepts
too few inputs to cover the slice's positive set. Quantitative
sparsification permits width $b=\Theta(\ln n/\ln\ln n)$ and hence a circuit
exponent of order $\sqrt{bn}$.

\section{Correlation and the two structural tools}\label{sec:tools}
The cube $\bits^V$ on a finite coordinate set $V$ has uniform measure.
Boolean functions are identified with $0/1$ indicators. A clause is an OR
of literals, and its width is the number of distinct variables after
normalization. A false clause has width zero; an empty CNF is true.
For $b\ge1$ and a real function $G$ on this cube, define
\[
 \corr_b(G)=\max_J |\E[GJ]|,
\]
where $J$ ranges over all width-at-most-$b$ CNF indicators, without a
clause-count restriction. Constants are allowed. There are finitely many
such indicators, so the maximum exists.

A $p$-random restriction retains each coordinate independently with
probability $p$, giving a live set $T$, and assigns the other coordinates
uniformly, giving $\rho$. Write $G_{T,\rho}$ for the remaining function.
Uniform completion of a random restriction is uniform on the full cube;
therefore
\begin{equation}\label{eq:restriction-average}
 \E_xG(x)H(x)=\E_{T,\rho}\E_zG_{T,\rho}(z)H_{T,\rho}(z).
\end{equation}

\begin{lemma}[Quantitative disjoint sparsification]\label{lem:sparsification}
There exist an absolute constant $K\ge2$ and positive integers $M(b)$,
for $b\ge2$, such that every width-$b$ CNF indicator on $m\ge1$ variables
is a sum of at most $2^{m/8}$ pairwise disjoint width-$b$ CNF indicators,
each having at most $M(b)m$ clauses, with
\begin{equation}\label{eq:sparse}
 M(b)\le(Kb)^{Kb}.
\end{equation}
\end{lemma}
\begin{proof}[Derivation from quantitative sparsification]
We use the quantitative counting form stated in~\cite[Section~4,
p.~6:9]{Williams}, combining~\cite[Section~6]{CIP} with the disjoint
refinement in~\cite[Appendix~A, Lemma~A.1 and its following proof]{DHMTW}.
Fix an absolute error $\eta_0\in(0,1/8]$ in that bound. The quantitative
estimate $(b/\eta_0)^{O(b)}$ means that there are absolute constants
$A,C\ge1$, independent of both $b$ and $m$, for which the density before
adding branch-defining units is at most
\[
 D(b)\le A(b/\eta_0)^{Cb}.
\]
The finite sparsification recursion has at most $2^{\eta_0m}$ leaves
for every input on $m$ variables. Lemma~A.1 of~\cite{DHMTW} states these
finite quantifiers explicitly; they are bounds on the leaves and formulas,
not just an asymptotic running-time estimate.

Here is the disjointness modification. At a flower with heart clause
$H=\bigvee_{\ell\in U}\ell$ and petals $P_i$, the relevant clauses are
$H\vee P_i$. The heart branch enforces $H$; the petal branch enforces
all $P_i$ and the units $\neg\ell$ for $\ell\in U$. These two branches
partition the satisfying assignments according to the truth of $H$.
Iterating gives disjoint leaves. A consistent leaf contains at most one
distinct added unit per variable, hence at most $m$; inconsistent leaves
are discarded. Width cannot increase. The reduction step removes clauses
properly containing a unit, so these units create no new flowers and the
leaf bound is unchanged, as proved in~\cite[Appendix~A]{DHMTW}.

Set $M(b)=\lceil D(b)\rceil+1$. For example, choosing any absolute
$K\ge\max\{2,2C,\eta_0^{-1},3A\}$ gives
\[
 M(b)\le D(b)+2\le3A(b/\eta_0)^{Cb}\le(Kb)^{Kb}
 \qquad(b\ge2).
\]
Also $2^{\eta_0m}\le2^{m/8}$. Both estimates hold for every finite $m$,
including when $b>m$. The identically false formula may be represented
by no nonzero leaves. This proves the claimed indicator identity.
No numerical values of $A,C,\eta_0$, or $K$ are needed or asserted here.
\end{proof}

\begin{lemma}[Restriction bound from~\cite{OAI26}]\label{lem:restriction}
Let $k\ge0$ and $b\ge1$ be integers, let $H$ be a width-$k$ CNF on $d$
variables, and let $0<p<1$. If
\begin{equation}\label{eq:theta}
 \theta=\sum_{\ell=b+1}^k\binom{k}{\ell}
             \left(\frac{2p}{1-p}\right)^\ell<1,
\end{equation}
there is a nonnegative $W_H=W(H,T,\rho)$ such that
\[
 \E W_H\le(1-\theta)^{-1},\qquad
 |\E_z G(z)H_{T,\rho}(z)|\le W_H\corr_b(G)
\]
simultaneously for every real function $G$ on the live cube.
\end{lemma}
The full finite proof appears in Appendix~\ref{app:restriction}. It applies
to each finite choice of $b,k,d,p$; no fixed-width limiting argument is
part of the lemma.

\section{The language and finite hard-slice estimates}\label{sec:language}
We recall the evaluator from~\cite[Section~4]{OAI26}. Its use of parameter
bits to select hard subfunctions is related to the input-index method
of Paturi, Saks, and Zane~\cite[Section~5]{PSZ00}.

\begin{definition}[The fixed language]\label{def:language}
For input length $n$, set $d=\lfloor n/5\rfloor$ and reject if $d<64$.
Otherwise set
\[
 r=\lceil d^{2/3}\rceil,\qquad
 t=\max\{j\ge0:j\text{ is even and }j^6\le d\}.
\]
As a safety check, reject if fewer than $2d+(t+2)r-1$ bits are available;
Proposition~\ref{prop:polytime} shows this branch is never taken for $d\ge64$.
Parse the input,
in order, into $d$ data bits $x$, $d+r-1$ hash bits $u$, $r$ modulus bits
$p_0,\ldots,p_{r-1}$, and $t$ coefficient blocks of $r$ bits each.
Ignore unused bits. The modulus bits define
\[
 P(Z)=Z^r+\sum_{i=0}^{r-1}p_iZ^i,
 \qquad R_P=\mathbb F_2[Z]/(P),
\]
and the coefficient blocks define $\beta_0,\ldots,\beta_{t-1}\in R_P$
in the basis $1,Z,\ldots,Z^{r-1}$.
Let
\begin{equation}\label{eq:hash}
 h_i(x)=\sum_{a=1}^d u_{i+a-1}x_a\pmod2\quad(1\le i\le r),
 \qquad h_u(x)=\sum_{i=1}^rh_i(x)Z^{i-1}.
\end{equation}
If $\lambda:R_P\to\mathbb F_2$ selects the constant coefficient of the
representative of degree less than $r$, accept exactly when
\begin{equation}\label{eq:evaluator}
 \lambda\!\left(\sum_{j=0}^{t-1}\beta_jh_u(x)^j\right)=0.
\end{equation}
\end{definition}

\begin{proposition}\label{prop:polytime}
Definition~\ref{def:language} gives one deterministic polynomial-time
language. For all sufficiently large $d$,
\begin{equation}\label{eq:rt}
 r\ge d^{2/3},\qquad \tfrac12d^{1/6}\le t\le d^{1/6}.
\end{equation}
\end{proposition}
\begin{proof}
For $d\ge64$, $r\le2d^{2/3}$ and $t+2\le2d^{1/6}$. Thus the blocks use
$2d+(t+2)r-1\le2d+4d^{5/6}-1\le4d-1\le n$ bits.
Also $d^{1/6}-2<t\le d^{1/6}$, which gives~\eqref{eq:rt} for
$d\ge4096$. Already $t\ge2$ for $d\ge64$, which is all that the finite
criterion below requires. The integer parameters can be found by testing $r^3\ge d^2$
and the sixth powers of even integers, with polynomial overhead.

Every monic modulus gives a well-defined quotient ring, whether or not it
is irreducible. Schoolbook multiplication and monic reduction cost
$O(r^2)$ bit-array operations; the hash costs $O(rd)$; and Horner evaluation
costs $O(tr^2)$. These fixed loops are implementable by one Turing machine
with polynomial overhead. The evaluator neither searches for hard
parameters nor tests irreducibility.
\end{proof}

\begin{lemma}[Finite hard-slice criterion]\label{lem:slices}
Let $d\ge64$, with $r,t$ as above, and let $b\ge2$, $0<p<1$, and
$\mu=pd\ge2$. Set $M=M(b)$ from Lemma~\ref{lem:sparsification}.
Assume $r\ge(33/16)\mu$ and, for every integer $m\in[\mu/2,2\mu]$,
\begin{equation}\label{eq:union}
 \frac{M[1+b\log_2(2m+1)]}{t}+\frac{\log_2t}{m}\le\frac18.
\end{equation}
There is a fixed setting of the non-data bits whose sign
$g:\bits^d\to\{-1,1\}$, with $+1$ on acceptance, satisfies
\begin{equation}\label{eq:finite-good}
 \Prb_{T,\rho}\!\left[
 |T|\notin[\mu/2,2\mu]\ \text{or}\
 \corr_b(g_{T,\rho})>2^{-|T|/8}\right]
 \le4\cdot2^{-\mu/16}.
\end{equation}
\end{lemma}
\begin{proof}
Choose a monic irreducible $P$ of degree $r$, which exists over
$\mathbb F_2$ for each $r\ge1$, and identify $R_P$ with
$F=\mathbb F_{2^r}$. Choose $u$ and the coefficients $\beta_j$ independently
and uniformly. We first bound the joint failure probability over these
choices and the random restriction.

\paragraph{Hash injectivity.}
For nonzero $w\in\bits^d$, let $a$ be its largest nonzero coordinate.
The $i$th output bit of $h_u(w)$ contains $u_{i+a-1}$ with coefficient one,
and all other indices in that row are smaller. These successive pivots
show that the linear map $u\mapsto h_u(w)$ has rank $r$. Hence its image
is uniform and $\Prb_u[h_u(w)=0]=2^{-r}$. A restricted data cube of dimension
$m$ has $2^m-1$ nonzero difference vectors, all supported on $T$.
Linearity and a union bound give
\begin{equation}\label{eq:collision}
 \Prb_u[h_u\text{ is not injective on this cube}]
 \le(2^m-1)2^{-r}\le2^{m-r}.
\end{equation}

\paragraph{Limited independence.}
Fix an injective hash on the cube. For any $q\le t$ distinct points, the
hashes are distinct elements $a_1,\ldots,a_q\in F$. Lagrange interpolation
shows that evaluation of degree-less-than-$t$ polynomials onto $F^q$ is
surjective. Uniform coefficients therefore give independent uniform
values at those points. Since $\lambda$ is a nonzero binary linear
functional, applying $(-1)^\lambda$ yields unbiased $t$-wise independent
signs. Thus the following moment bound is conditional on any injective hash.

\paragraph{Sparse tests.}
There are at most $(2m+1)^b$ normalized clauses of width at most $b$:
encode their literals in a fixed order and pad to length $b$. Include the
empty false clause. Lists of $Mm$ clauses, with an extra padding symbol,
then encode at most
\begin{equation}\label{eq:test-count}
 2^{Mm[1+b\log_2(2m+1)]}
\end{equation}
CNF indicators. This includes both constants.

For one such indicator $J$, put $Y=\sum_{z\in\bits^T}g_{T,\rho}(z)J(z)$.
In the expansion of $Y^t$, a tuple with a point occurring exactly once has
zero expectation. Every remaining tuple uses at most $t/2$ distinct points.
Listing $t/2$ points, with padding, and a choice of list position for each
factor gives at most $t^t2^{mt/2}$ tuples. Since $t$ is even,
\[
 \E|Y|^t\le t^t2^{mt/2},\qquad
 \Prb[|\E_z g_{T,\rho}J|>2^{-m/4}]\le t^t2^{-mt/4}.
\]
Combining this with~\eqref{eq:test-count}, the probability that any sparse
test fails is at most
\[
 2^{Mm[1+b\log_2(2m+1)]+t\log_2t-mt/4}
 \le2^{mt(1/8-1/4)}=2^{-mt/8},
\]
where the inequality is exactly~\eqref{eq:union}. On the complementary
event, Lemma~\ref{lem:sparsification} expresses every width-$b$ CNF as a
sum of at most $2^{m/8}$ tests from this same sparse family. The triangle
inequality gives $\corr_b(g_{T,\rho})\le2^{-m/8}$ simultaneously for all
width-$b$ CNFs.

\paragraph{Restriction tails and averaging.}
The live dimension $m$ is binomial with mean $\mu$. Independence gives
\[
 \E2^{-m}=(1-p/2)^d\le e^{-\mu/2},\qquad
 \E2^m=(1+p)^d\le e^\mu.
\]
Markov's inequality yields
\begin{align*}
 \Prb[m<\mu/2]&\le2^{-((\log_2e-1)/2)\mu},\\
 \Prb[m>2\mu]&\le2^{-(2-\log_2e)\mu}.
\end{align*}
Both exponent coefficients exceed $1/16$, so their sum is at most
$2\cdot2^{-\mu/16}$. Inside the dimension window,
$r\ge(33/16)\mu$ gives $2^{m-r}\le2^{-\mu/16}$, while $t\ge2$ gives
$2^{-mt/8}\le2^{-\mu/16}$. Adding the two errors and the two dimension
tails bounds joint failure by $4\cdot2^{-\mu/16}$.
Averaging over the finitely many settings of $(u,\beta_0,\ldots,\beta_{t-1})$
leaves one setting with at most this failure probability over restrictions.
Fix it, the chosen $P$, and arbitrary unused bits.
\end{proof}

\section{Growing parameters and the lower bound}\label{sec:main-proof}


We now prove Theorem~\ref{thm:main}. All unqualified logarithms are natural.
For sufficiently large $d$, set
\begin{align*}
 b&=\left\lfloor\frac{\ln d}{24K\ln\ln d}\right\rfloor,&
 s&=\frac{\sqrt b}{256},& B&=64(s+1),\\
 p&=\frac{B}{\sqrt d},& \mu&=B\sqrt d,&
 k&=\left\lceil3(s+1)\sqrt d\right\rceil.
\end{align*}
We work beyond a fixed absolute threshold so that $b\ge65536$, $s\ge1$,
and $p\le1/2$. Let $M=M(b)$ from Lemma~\ref{lem:sparsification}.

\paragraph{The sparse-test union bound.}
Eventually $\ln(Kb)\le2\ln\ln d$, so \eqref{eq:sparse} gives
\[
 \ln M\le Kb\ln(Kb)\le\frac{\ln d}{12},\qquad M\le d^{1/12}.
\]
Uniformly for $\mu/2\le m\le2\mu$, we have
$\log_2(2m+1)=O(\ln d)$ and $b=O(\ln d)$.
More explicitly, \eqref{eq:rt} gives the finite envelope
\begin{align*}
 \frac{M[1+b\log_2(2m+1)]}{t}+\frac{\log_2t}{m}
 &\le 2M d^{-1/6}
       [1+b\log_2((4B+1)\sqrt d)]
       +\frac{\ln d}{3B\ln2\,\sqrt d}\\
 &\le 2d^{-1/12}
       [1+b\log_2((4B+1)\sqrt d)]
       +\frac{\ln d}{3B\ln2\,\sqrt d}.
\end{align*}
Consequently,
\[
 \frac{M[1+b\log_2(2m+1)]}{t}
       +\frac{\log_2t}{m}
 =O\!\left(d^{-1/12}(\ln d)^2
                   +\frac{\ln d}{\sqrt d}\right)=o(1).
\]
The constants here are absolute once the absolute $K$ is fixed.
This proves \eqref{eq:union} uniformly over the entire window.

\paragraph{Hash and restriction tails.}
Since $B=O(\sqrt{\ln d})$, we have $r/\mu\to\infty$.
In particular, eventually $r\ge(33/16)\mu$. Within the window,
\[
 2^{m-r}\le2^{-\mu/16},\qquad
 2^{-mt/8}\le2^{-\mu/16}.
\]
For the binomial random variable $m=|T|$ of mean $\mu$, the
exponential-moment bounds in Lemma~\ref{lem:slices} give
\[
 \Prb[m\notin[\mu/2,2\mu]]\le2\cdot2^{-\mu/16}.
\]
Averaging the joint failure probability over the parameter bits therefore
gives one setting of those bits, for this length, whose sign function $g$
satisfies
\begin{equation}\label{eq:good}
 \Prb\left[m\notin[\mu/2,2\mu]
       \text{ or }\corr_b(g_{T,\rho})>2^{-m/8}\right]
       \le4\cdot2^{-\mu/16}.
\end{equation}
This is a pointwise existence argument at each $d$; it does not exchange
the original fixed-parameter asymptotic quantifiers.

\paragraph{The restriction lemma still applies.}
For $d\ge1$,
\[
 k\le[3(s+1)+1]\sqrt d,\qquad
 \frac{2pk}{1-p}\le4B[3(s+1)+1]=:Q.
\]
Using $B=64(s+1)$ and $s\ge1$,
\[
 Q\le1024(s+1)^2\le4096s^2=\frac b{16}.
\]
By $\binom{k}{\ell}\le k^\ell/\ell!$ and
$\ell!\ge(\ell/e)^\ell$,
\[
 \theta\le\sum_{\ell>b}\frac{Q^\ell}{\ell!}
 \le\sum_{\ell>b}(e/16)^\ell<\frac12.
\]
Thus Lemma~\ref{lem:restriction} supplies $\E W_H\le2$ for every width-$k$ test $H$.

\subsection{Correlation transfer and the circuit contradiction}

On the good restrictions in \eqref{eq:good}, the restricted correlation
is at most $W_H2^{-\mu/16}$. On the bad restrictions it is at most one.
Uniform completion of a random restriction is uniform on $\bits^d$.
Taking expectations and using the triangle inequality gives
\begin{equation}\label{eq:corr}
 |\E_x g(x)H(x)|\le6\cdot2^{-\mu/16}
\end{equation}
for every width-$k$ CNF $H$. Independence of $W_H$ and the good event
is not required. Applying \eqref{eq:corr} to the constant-true CNF
shows that $\Prb[g=1]\ge1/4$ for sufficiently large $d$.

Suppose there were a circuit for $f_n$ with at most $S\le2^{s\sqrt d}$
total gates. Fix its parameter inputs to the setting chosen above.
Delete gates not connected to the output. Each remaining middle AND gate
is a CNF $H$ with at most $S$ distinct bottom clauses, and there are at
most $S$ middle gates, even with shared bottom gates. Normalize its clauses
by deleting true clauses and repeated literals; tautological clauses are
true and are removed. A middle gate with a false clause computes zero
and may be omitted.
Delete from $H$ the clauses of width greater than $k$, obtaining $H'$.
Then $H'\ge H$ and
\[
 \E(H'-H)\le S2^{-k}.
\]
Every input accepted by $H$ is accepted by the top OR, so $gH=H$.
Equation \eqref{eq:corr} therefore implies
\[
 \E H\le6\cdot2^{-\mu/16}+S2^{-k}.
\]
Taking the union over middle gates yields
\begin{align*}
 \Prb[g=1]
 &\le6S2^{-\mu/16}+S^2 2^{-k}\\
 &\le6\cdot2^{-(3s+4)\sqrt d}
             +2^{-(s+3)\sqrt d}=o(1),
\end{align*}
contradicting $\Prb[g=1]\ge1/4$.
Hence $S_3(f_n)>2^{s\sqrt d}$ for every sufficiently large $n$.
Finally,
\[
s\sqrt d=\frac{\sqrt{bd}}{256}
       =\Theta\!\left(\sqrt{\frac{n\ln n}{\ln\ln n}}\right),
\]
because $d=\lfloor n/5\rfloor$ and
$b=\Theta(\ln d/\ln\ln d)$. More explicitly, eventually
\[
 d\ge\frac n6,\qquad \ln d\ge\frac{\ln n}{2},\qquad
 b\ge\frac{\ln d}{48K\ln\ln d}
   \ge\frac{\ln n}{96K\ln\ln n}.
\]
It follows that
\[
 s\sqrt d\ge\frac1{6144\sqrt K}
                 \sqrt{\frac{n\ln n}{\ln\ln n}},
\]
which proves Theorem~\ref{thm:main}. \hfill$\square$


\section{Remarks on scope}
The gain in the exponent is the unbounded factor
$\sqrt{\ln n/\ln\ln n}$. It remains smaller than every fixed positive
power of $n$, so this argument does not give $2^{n^{1/2+\varepsilon}}$ for
fixed $\varepsilon>0$. The role of quantitative sparsification is specific:
$M(b)\le\exp(O(b\log b))$ remains small enough compared with the
$d^{1/6}$ independence parameter when $b$ is of order $\ln d/\ln\ln d$.
The relation $s=\sqrt b/256$ then converts this width into the asserted
circuit exponent.

\paragraph{Size of the threshold.}
The proof is asymptotic, with extremely conservative constants. The
requirement $b\ge65536$ alone implies
\[
 \frac{\ln d}{\ln\ln d}\ge24K\cdot65536\ge3{,}145{,}728.
\]
On the large-$d$ branch this requires $d$ to have more than 24 million
decimal digits. The other eventual inequalities may increase the threshold
further. This is a limitation of the displayed parameter schedule, not a
claim about the smallest lengths at which the lower bound is true.

\paragraph{The logarithmic loss.}
With the available estimate $M(b)\le\exp(O(b\ln b))$, certifying that
$M(b)$ stays within a fixed power of $d$ leads to
$b\ln b=O(\ln d)$. This accounts for the $\ln\ln n$ loss. A hypothetical
uniform disjoint sparsification theorem with $M(b)\le2^{Cb}$ for an
absolute $C$ and the same $2^{m/8}$ piece bound would permit
$b=\delta\ln d$ for a sufficiently small absolute $\delta>0$.
The same estimates would then give $2^{\Omega(\sqrt{n\ln n})}$ for this
language, without implying the conjectural majority bound.
Section~7 of~\cite{CIP} asks about a corresponding improvement in
sparsification density. We do not assert that such a disjoint improvement
is impossible; the present density estimate supplies a sufficient bound,
not a matching lower bound.

Increasing $t$ to another fixed power of $d$, while maintaining the input
budget and hash estimates, would enlarge the polynomial density allowance.
With the same sparsification estimate it would change constants, not the
order $\ln d/\ln\ln d$ of the certified width. Moreover, $t$ determines the
number of coefficient blocks, so changing it would define a variant of
the language rather than merely retune its proof.

All choices of parameter bits are existential restrictions of one fixed
language. They may depend on $n$ and on the proof parameters, while the
membership algorithm does not. The proof uses no uniformity assumption
on the competing circuits.

\appendix
\section{The finite restriction argument}\label{app:restriction}
For completeness, we reproduce the argument of~\cite[Section~2]{OAI26}
in the notation of Lemma~\ref{lem:restriction}. The two points to check are
an exact finite signed expansion into width-$b$ tests and a bound on its
expected absolute coefficient mass. Counting paths backwards is what
removes the number of clauses from that bound.

\begin{proof}[Proof of Lemma~\ref{lem:restriction}]
Normalize $H$ by deleting tautological and true clauses and deleting repeated
literals within each remaining clause. Keep clause indices, including
repeated clauses. Write $E_i$ for the scope of clause $C_i$. A violation of
$C_i$ specifies a unique pattern on $E_i$; the false clause has empty scope.
Fix a restriction $(T,\rho)$, and let $v_i(z)$ be its violation indicator.
For $U\subseteq T$ set
\[
 q=\sum_i v_i,\qquad
 a_U=\sum_{i:\,|E_i\cap(T\setminus U)|\le b}v_i,\qquad
 c_U=\sum_{i:\,E_i\cap(T\setminus U)=\varnothing}v_i.
\]
Call an index hard at $U$ if $|E_i\cap(T\setminus U)|>b$.
The function $K_0=\mathbf1[a_\varnothing=0]$ is a width-$b$ CNF.

A path $P=(i_1,\ldots,i_j)$ starts at $U_0=\varnothing$, selects an index
hard at $U_{h-1}$, and sets $U_h=U_{h-1}\cup(E_{i_h}\cap T)$.
Each step adds at least $b+1$ coordinates, so paths have length at most
$L=\lfloor |T|/(b+1)\rfloor$. On the support of
$\prod_{h=1}^j v_{i_h}$, the selected indices are distinct and
$a_{U_h}\ge c_{U_h}\ge h$.
All fractions below are defined as zero when their numerator vanishes.

For a nonempty path define
\[
 r_P=\frac{\prod_{h=1}^jv_{i_h}}
                {q\prod_{h=1}^{j-1}a_{U_h}},\qquad
 t_P=\frac{\prod_{h=1}^jv_{i_h}}
                {\prod_{h=1}^{j}a_{U_h}}.
\]
Since $q=a_{U_j}+\sum_{i\text{ hard at }U_j}v_i$,
\[
 r_P=t_P-\sum_{i\text{ hard at }U_j}r_{Pi}.
\]
At a point with $K_0=1$ and $q>0$, all violated clauses are initially hard,
and $\sum_{|P|=1}r_P=1$. Iterating the displayed identity until there are
no further paths gives
\begin{equation}\label{eq:audited-expansion}
 H_{T,\rho}=K_0+
   \sum_{j=1}^{L}(-1)^j\sum_{|P|=j}K_0t_P.
\end{equation}
At points with $K_0=0$ every term vanishes; at points with $q=0$ the
nonempty-path terms vanish and $K_0=H_{T,\rho}=1$.
Thus this is an identity on the whole cube, not just on its error set.

Call a path compatible with $\rho$ when its selected violating patterns
are mutually consistent and agree with $\rho$ on fixed coordinates.
For such a path, all selected violations prescribe a single pattern
$\sigma_P$ on $U_j$. Denote its cylinder by $Z_P$.
The numerator of $t_P$ is $\mathbf1_{Z_P}$.
On $Z_P$ the function $c_{U_h}$ is a positive constant $c_h$, because all
variables of every counted clause have been fixed. Moreover,
\[
 a_{U_h}=c_h+w_h,
 \qquad
 w_h=\sum_{i:\,0<|E_i\cap(T\setminus U_h)|\le b}v_i.
\]
After substituting the final pattern $\sigma_P$, the indexed clauses
counted by $w_h$ have width at most $b$ on the remaining coordinates.
Their multiplicities and any constant clauses are retained.

For any indexed family of clauses and positive integer $c$, randomly
permute its clause objects together with $c$ distinct markers, and take
the conjunction of the clauses preceding the first marker. On an assignment
violating $w$ of those clauses, this conjunction holds with probability
$c/(c+w)$: the earliest of the $w$ violated clause objects and $c$ markers
must be a marker. This is a finite convex mixture of CNF indicators.
Apply the construction independently for each $h$, then conjoin the sampled
CNFs with $K_0$ and the unit clauses defining $Z_P$. Every resulting test
$J_P$ has width at most $b$, using $b\ge1$. It satisfies
\[
 K_0t_P=\omega(P)\,\E J_P,
 \qquad \omega(P)=\prod_{h=1}^j\frac1{c_h}.
\]
An incompatible path contributes zero. Define
\[
 W_0=1,\qquad
 W_j=\sum_{\substack{|P|=j\\P\text{ compatible with }\rho}}\omega(P),
 \qquad W=\sum_{j=0}^{\lfloor d/(b+1)\rfloor}W_j.
\]
Equation~\eqref{eq:audited-expansion} and the triangle inequality prove
$|\E_zGH_{T,\rho}|\le W\corr_b(G)$ for every $G$ simultaneously.
The construction of $W$ never uses $G$.

It remains to prove $\E W_j\le\theta^j$ for $j\ge1$.
Generate a uniform full assignment $y$, independently of $T$, and set
$\rho=y|_{[d]\setminus T}$. Conditional on $(T,\rho)$, a compatible path
has all its clauses violated by $y$ with probability $2^{-|U_j|}$.
Consequently
\begin{equation}\label{eq:audited-reweight}
 \E W_j=\E_{T,y}
 \sum_{\substack{|P|=j\\C_{i_1},\ldots,C_{i_j}\text{ violated by }y}}
 2^{|U_j|}\prod_{h=1}^j\frac1{c_{U_h}(y|_T)}.
\end{equation}
This is conditional cylinder averaging; it makes no claim of independence
between the path and its denominators.

Fix $y$, and for a remainder $R\subseteq[d]$ put
\[
 V_R=\{i:C_i(y)=0,\ E_i\cap R=\varnothing\}.
\]
An actual path has $R_h=T\setminus U_h$ and
$D_h=U_h\setminus U_{h-1}$, with
\[
 i_h\in V_{R_h},\qquad D_h\subseteq E_{i_h},\qquad |D_h|>b,
 \qquad R_{h-1}=R_h\mathbin{\dot\cup}D_h,\qquad
 c_{U_h}(y|_T)=|V_{R_h}|.
\]
Its terminal remainder $R_j$ and the reverse sequence
$(i_j,D_j),\ldots,(i_1,D_1)$ determine both $T$ and its entire indexed
path. Hence this is an injective encoding even when clauses repeat.

Let $\alpha=2p/(1-p)$ and
$\pi(R)=p^{|R|}(1-p)^{d-|R|}$. The weight of an encoded path in
\eqref{eq:audited-reweight}, including its live-set probability, is exactly
\[
 \pi(T)2^{|U_j|}\prod_h\frac1{|V_{R_h}|}
   =\pi(R_j)\prod_h\frac{\alpha^{|D_h|}}{|V_{R_h}|}.
\]
Enlarge the reverse sequences by allowing, at any state $R$, every choice
$i\in V_R$ and $D\subseteq E_i$ with $|D|>b$, followed by $R\cup D$.
Here $D\cap R=\varnothing$ because $i\in V_R$.
All weights are nonnegative, and the outgoing weight at each state with
$V_R\ne\varnothing$ is at most
\[
 \frac1{|V_R|}\sum_{i\in V_R}
          \sum_{\ell=b+1}^{|E_i|}\binom{|E_i|}{\ell}\alpha^\ell
 \le\sum_{\ell=b+1}^k\binom{k}{\ell}\alpha^\ell=\theta.
\]
If $V_R$ is empty there are no outgoing steps. Induction on the number
of reverse steps bounds their total weight from any state by $\theta^j$.
Summing over terminal remainders costs only
$\sum_R\pi(R)=1$. This proves $\E_T[\text{path mass}]\le\theta^j$ for every fixed $y$, and
therefore $\E W_j\le\theta^j$.

Finally $\E W\le\sum_{j=0}^{\lfloor d/(b+1)\rfloor}\theta^j
\le(1-\theta)^{-1}$. Empty formulas, false clauses, repeated clauses,
empty live sets, and $k\le b$ all satisfy the same argument.
\end{proof}


\section{Verification and provenance}\label{app:verification}
The source construction and restriction argument were read at repository
revision \href{https://github.com/openai/math/commit/adc7f1241b42e322a6451854ab7e4b4c146bf78a}{\texttt{adc7f1241b42e}}. The accompanying
research note records the detailed literature comparison and search scope.
The manuscript date of the source is September 23, 2026; its inspected Git
revision is dated October 6, 2026. These dates alone do not establish
publication priority for an unstated consequence.

Two Python programs accompany the derivation. The file
\path{check_quantitative_depth_three.py} checks the elementary constants
with exact arithmetic and nine illustrative parameter choices in the log
domain. The chosen sample values of $K$ are not certified sparsification
constants. The file \path{check_restriction_lemma.py} checks all 512
sets of normalized clauses on two variables, and 11 additional formulas
at two rational retention probabilities. Its 534 formula--parameter cases
cover 11,254 restrictions and 43,524 completions, including nonzero
three-step paths and marker mixtures with seven indexed clauses.
All checks passed. These finite diagnostics support the proof but do not
replace its universal arguments.

The repository contains a Lean theorem for the restriction lemma and
finite numerical estimates. The strengthened rate and growing schedule
in this paper have not been integrated into that formalization, and no
fresh Lean kernel verification was performed for this work. The manuscript
was prepared with assistance from Codex, including derivation review,
literature search, computational diagnostics, and typesetting. It is a
research preprint and has not undergone external peer review.

\begin{thebibliography}{9}
\bibitem{OAI26}
OpenAI.
\emph{Beyond the Square-Root Exponent for Depth-Three Boolean Circuits}.
Manuscript dated September 23, 2026. Repository revision
\texttt{adc7f1241b42e322a6451854ab7e4b4c146bf78a}.
\href{https://github.com/openai/math/tree/adc7f1241b42e322a6451854ab7e4b4c146bf78a/preprints/Beyond-the-Square-Root-Exponent-for-Depth-Three-Boolean-Circuits-September-23-2026}{Source manuscript and LaTeX}.

\bibitem{CIP}
Chris Calabro, Russell Impagliazzo, and Ramamohan Paturi.
A duality between clause width and clause density for SAT.
In \emph{21st IEEE Conference on Computational Complexity},
pages 252--260, 2006. Section 6.
\href{https://doi.org/10.1109/CCC.2006.6}{doi:10.1109/CCC.2006.6}.
\href{https://cseweb.ucsd.edu/~ccalabro/cip.pdf}{Author manuscript}.

\bibitem{DHMTW}
Holger Dell, Thore Husfeldt, D\'aniel Marx, Nina Taslaman, and Martin Wahl\'en.
Exponential time complexity of the permanent and the Tutte polynomial.
\emph{ACM Transactions on Algorithms}, 10(4), article 21, 2014.
Appendix A of the \href{https://arxiv.org/abs/1206.1775v1}{2012 full version}.
\href{https://doi.org/10.1145/2635812}{doi:10.1145/2635812}.

\bibitem{Williams}
R. Ryan Williams.
Counting solutions to polynomial systems via reductions.
In \emph{SOSA 2018}, OASIcs 61, article 6, 2018. Section 4, p.~6:9.
\href{https://people.csail.mit.edu/rrw/counting-sosa.pdf}{Author manuscript}.

\bibitem{PSZ00}
Ramamohan Paturi, Michael E. Saks, and Francis Zane.
Exponential lower bounds for depth three Boolean circuits.
\emph{Computational Complexity}, 9(1):1--15, 2000.
\href{https://doi.org/10.1007/PL00001598}{doi:10.1007/PL00001598}.

\bibitem{GKKPRST}
Mohit Gurumukhani, Daniel Kleber, Ramamohan Paturi, Christopher Rosin,
Michael Saks, and Navid Talebanfard.
\emph{Optimal Monotone Depth-Three Circuit Lower Bounds for Majority}.
arXiv:2601.04072v2, February 3, 2026.
\href{https://arxiv.org/abs/2601.04072v2}{arXiv:2601.04072v2}.

\bibitem{GKPRT26IP}
Mohit Gurumukhani, Daniel Kleber, Ramamohan Paturi, Christopher Rosin,
and Navid Talebanfard.
\emph{Optimal Depth-Three Circuits for Inner Product}.
arXiv:2601.04446v2, February 12, 2026.
\href{https://arxiv.org/abs/2601.04446v2}{arXiv:2601.04446v2}.
\end{thebibliography}
\end{document}
